\documentclass{jpsj2}
\usepackage{graphicx}
% BibTeX's styles use \newblock, which the class doesn't define.
\providecommand{\newblock}{}

\title{Title of the Paper}

\author{First \textsc{Author}$^{1}$\thanks{first@example.org} and Second \textsc{Author}$^{2}$}
\inst{$^{1}$Department of Physics, University, City 000-0000, Japan\\
$^{2}$Institute, City 000-0000, Japan}

\abst{State the problem, what you did, and the main result in a few sentences.}

\kword{first keyword, second keyword, third keyword}

\begin{document}

\maketitle

\section{Introduction}
Introduce the problem and why it matters, then the contributions of this
paper.\cite{bardeen1957} Several works at once.\cite{bardeen1957,feynman1982}

\section{Method}
Describe what you did. Numbered equations use the \texttt{equation} environment:
\begin{equation}
  H = -t \sum_{\langle i,j \rangle, \sigma} c^\dagger_{i\sigma} c_{j\sigma} + U \sum_i n_{i\uparrow} n_{i\downarrow}.
  \label{eq:main}
\end{equation}

\section{Results}
Figure~\ref{fig:overview} shows the main result.

\begin{figure}
  \centering
  \includegraphics[width=\columnwidth]{example-image}
  \caption{What the figure shows. Upload an image and use its file name here.}
  \label{fig:overview}
\end{figure}

\section{Conclusions}
Summarise the findings and what comes next.

\begin{acknowledgment}
Thank the people and funding that supported this work.
\end{acknowledgment}

\bibliographystyle{unsrt}
\bibliography{references}

\end{document}
